\documentclass[UTF8,a4paper,11pt]{ctexart}

\usepackage[a4paper,margin=2.25cm,headheight=15pt]{geometry}
\usepackage{amsmath,amssymb}
\usepackage{booktabs,longtable,tabularx,array}
\usepackage{enumitem}
\usepackage{xcolor}
\usepackage{listings}
\usepackage{fancyhdr}
\usepackage{lastpage}
\usepackage{hyperref}

\definecolor{accent}{HTML}{1F5A94}
\definecolor{softblue}{HTML}{EEF5FB}
\definecolor{softgray}{HTML}{F5F6F7}
\definecolor{codegreen}{HTML}{26734D}
\hypersetup{
  colorlinks=true,
  linkcolor=accent,
  urlcolor=accent,
  citecolor=accent,
  pdftitle={计图人工智能挑战赛正式赛算法与复现说明},
  pdfauthor={lc666}
}

\pagestyle{fancy}
\fancyhf{}
\lhead{计图人工智能挑战赛正式赛}
\rhead{基于图学习的推荐任务}
\cfoot{第 \thepage{} 页，共 \pageref{LastPage} 页}
\renewcommand{\headrulewidth}{0.4pt}

\setlist[itemize]{leftmargin=2em,itemsep=0.25em,topsep=0.35em}
\setlist[enumerate]{leftmargin=2.2em,itemsep=0.25em,topsep=0.35em}
\renewcommand{\arraystretch}{1.22}
\newcolumntype{Y}{>{\raggedright\arraybackslash}X}
\newcolumntype{C}[1]{>{\centering\arraybackslash}p{#1}}

\lstdefinestyle{codeblock}{
  backgroundcolor=\color{softgray},
  basicstyle=\ttfamily\small,
  breaklines=true,
  columns=fullflexible,
  frame=single,
  rulecolor=\color{black!15},
  keywordstyle=\color{accent},
  commentstyle=\color{codegreen},
  showstringspaces=false,
  tabsize=2
}
\lstset{style=codeblock}

\title{\vspace{-1.3cm}\textbf{计图人工智能挑战赛正式赛\\算法原理与完整复现说明}\\[0.7em]
\large 基于图学习的推荐任务（Track 1）}
\author{队伍：lc666}
\date{2026 年 8 月}

\begin{document}

\maketitle
\thispagestyle{empty}
\vfill
\begin{center}
\begin{minipage}{0.88\textwidth}
\colorbox{softblue}{%
  \parbox{0.94\textwidth}{
    \textbf{方案摘要：}将任务建模为时间有向图上的候选边排序。主分支以训练历史构造边频次、时效性、转移、跳跃转移、节点窗口和无标签候选分布特征，在正向图与反向图上提取时间加权低秩表示，再用 XGBRanker 学习每个查询行内部 100 个候选节点的相对次序。V69 主分支本地时间验证 MRR 为 1.3952170297；最终 V83 进一步使用数据集感知的双分支路由，dataset1 保留主分支，dataset2 选择互补预测分支，A 榜成绩为 1.4591。
  }}
\end{minipage}
\end{center}
\vfill
\begin{center}
\small
代码仓库：\url{https://github.com/purplellllll/jittor-lc666-graph-learning-recommendation}\\
GitLink：\url{https://gitlink.org.cn/llllllc/lc666}
\end{center}
\clearpage

\tableofcontents
\clearpage

\section{任务定义与建模}

\subsection{输入与输出}

训练集中的每条记录表示时间有向边
\[
e=(u,v,t),
\]
其中 $u$ 为源节点、$v$ 为目标节点、$t$ 为时间戳。测试集中的每个查询为
\[
q_i=(u_i,t_i,\mathcal{C}_i),\qquad
\mathcal{C}_i=\{c_{i1},c_{i2},\ldots,c_{i100}\}.
\]
模型需要为 100 个候选目标节点分别输出一个分数，真实目标节点的排序越靠前越好。最终结果按行输出 100 个浮点分数，列顺序与候选列 \texttt{c1--c100} 完全一致。

\subsection{评价指标}

本地验证使用平均倒数排名（Mean Reciprocal Rank, MRR）：
\[
\operatorname{MRR}=\frac{1}{N}\sum_{i=1}^{N}\frac{1}{r_i},
\]
其中 $r_i$ 是第 $i$ 个查询中真实目标的名次。该指标强烈奖励前几名，因此优化的重点不是估计绝对点击概率，而是提升\textbf{同一查询行内}真实节点的相对顺序。

\subsection{核心建模判断}

本任务同时包含三个结构：

\begin{itemize}
  \item \textbf{时间结构：}历史边越新，通常越能表示当前偏好；
  \item \textbf{有向图结构：}源节点、目标节点以及正反向关系的语义不对称；
  \item \textbf{候选集合结构：}模型只需在当前 100 个候选中排序，行内相对量比全局绝对量更稳定。
\end{itemize}

因此最终方案采用“规则图分数作为稳定先验 + 图与时间特征 + Learning-to-Rank 精排”的组合，而不是将问题处理成普通的独立二分类。

\section{整体算法流程}

完整流程可以概括为：

\begin{enumerate}
  \item 对训练边按时间稳定排序，构造历史图；
  \item 严格按时间切分训练历史和验证边，并使用正式测试候选行构造 100 候选验证查询；
  \item 从历史边构建 pair、transition、skip-transition、节点统计和最近序列索引；
  \item 在正向图与反向图上分别计算 32 维时间加权 TruncatedSVD；
  \item 为每个 $(u,t,c)$ 候选三元组生成 180 维原始特征，并按数据集选择特征子集；
  \item 以每个查询行为一个 group，训练目标为 \texttt{rank:ndcg} 的 XGBRanker；
  \item 将排序器分数与稳定规则基线按验证确定的系数混合；
  \item 使用完整训练历史重建全部索引，分块生成 V69 主分支的两个数据集预测；
  \item 校验第二预测分支的 ZIP 完整性、候选列数、行数、有限值与文件指纹；
  \item 根据数据集级验证结论执行分支路由：dataset1 取主分支，dataset2 取互补分支，再打包最终结果。
\end{enumerate}

用函数关系表示，最终候选分数为
\[
s_{i,j}=\alpha_d\,\widetilde{s}^{\text{ranker}}_{i,j}
+(1-\alpha_d)\,\widetilde{s}^{\text{base}}_{i,j},
\]
其中 $d$ 表示数据集，波浪号表示在查询行内执行 min--max 归一化。dataset1 的 $\alpha_1=0.95$，dataset2 的 $\alpha_2=1.00$。

上式得到主分支分数。最终提交分数为
\[
s^{\mathrm{final}}_{d,i,j}=
\begin{cases}
s^{\mathrm{graph}}_{i,j}, & d=\mathrm{dataset1},\\
s^{\mathrm{comp}}_{i,j}, & d=\mathrm{dataset2},
\end{cases}
\]
其中 $s^{\mathrm{comp}}$ 是互补预测分支对同一 100 候选集合给出的排序分数。最终选择在数据集层完成，不把两个尺度不同的原始分数直接相加。

\section{时间验证设计}

\subsection{严格历史切分}

验证必须模拟“用过去预测未来”。dataset1 按时间顺序使用前 85\% 边作为历史、后 15\% 边作为验证事件；dataset2 若官方训练数据提供 \texttt{split} 列，则使用 \texttt{split=0} 作为历史，\texttt{split=1} 作为验证。验证边数量最多采样 130,000 条，并继续保持时间顺序。

\subsection{候选行模拟}

对每条验证边 $(u,v,t)$，从正式测试集随机抽取一整行 100 候选作为负样本背景，再在随机位置插入真实目标 $v$。若候选行中已经包含 $v$，则替换重复项，使每行只有一个正样本。最终预设的训练候选模式和评估候选模式均为 \texttt{row}。

该策略保留了正式测试候选的联合分布与行内难度，比从全部候选节点独立采负样本更接近真实评测。随后再按时间顺序将模拟查询的前 90\% 用于排序器训练，后 10\% 用于本地评估。

\subsection{防止信息泄漏}

验证特征只由历史部分构建，验证真实目标不会参与历史计数、SVD 或转移统计。最终训练才使用完整训练集重建历史。测试集只读取源节点、时间和公开候选 ID，不读取、推断或构造测试真实标签。

\section{图统计与时间特征}

\subsection{直接边统计}

对历史有向边 $(u,v)$ 统计出现次数 $n(u,v)$ 和最后出现时间 $t_{\mathrm{last}}(u,v)$。使用
\[
f_{\mathrm{pair\_cnt}}=\log(1+n(u,v))
\]
压缩长尾计数，并用指数衰减表达时效性：
\[
f_{\mathrm{pair\_rec}}
=\exp\left(-\frac{\max(t-t_{\mathrm{last}}(u,v),0)}{\tau}\right),
\quad
\tau=\max\left(\frac{t_{\max}-t_{\min}}{10},1\right).
\]
同时计算反向 pair，即 $n(v,u)$ 及其新近程度，用于捕获双向关系或角色互换模式。

\subsection{节点热度与时间窗口}

对目标节点统计全局入边次数、最近出现时间，以及训练时间分位点 50\%、70\%、85\%、90\%、95\%、97\%、99\% 之后的窗口计数。源节点侧统计出边次数和最后活跃时间；候选节点也作为源节点统计其出边活跃程度。它们共同区分“长期热门”“短期爆发”和“近期沉寂”等状态。

\subsection{转移与跳跃转移}

对同一源节点 $u$ 的历史目标序列
\[
v_1,v_2,\ldots,v_L
\]
构造相邻转移 $v_{k-1}\rightarrow v_k$。为降低只使用一步转移的稀疏性，再统计最大间隔为 5 的跳跃转移：
\[
v_{k-\ell}\rightarrow v_k,\qquad \ell\in\{1,2,3,4,5\}.
\]
对查询源节点最近访问的 10 个目标分别查询正向/反向的相邻转移次数、跳跃转移次数和相应时间衰减值。这样既能描述短期序列偏好，也能容忍重复周期和非严格相邻行为。

\subsection{最近历史交互}

保存每个源节点最近 10 个目标，生成候选是否直接命中最近目标、匹配深度、加权匹配强度，以及源节点最近目标序列与候选节点最近序列的交集特征。交集数量和带位置衰减的交集强度可以近似表达局部共同邻居关系。

\section{时间加权双向图 SVD}

\subsection{正向低秩表示}

设图的稀疏邻接矩阵为 $A\in\mathbb{R}^{|V|\times|V|}$。每条历史边按其在全局时间跨度中的位置加权：
\[
w_e=1+0.35\cdot\frac{t_e-t_{\min}}{\max(t_{\max}-t_{\min},1)}.
\]
相同节点对的权重累加形成 $A_{uv}$。对 $A$ 做随机化 TruncatedSVD：
\[
A\approx U_k\Sigma_kV_k^{\top},\qquad k=32.
\]
实现中源节点表示为 \texttt{fit\_transform} 得到的 $U_k\Sigma_k$，目标节点表示为 $V_k$。对候选 $(u,c)$ 计算内积、余弦相似度、行内归一化值、倒数名次和两侧向量范数。

\subsection{反向低秩表示}

有向图中“谁指向谁”不可互换。方案在反向矩阵 $A^{\top}$ 上独立执行同样的 32 维分解，随机种子在主种子基础上加 17。正反向表示共同刻画节点作为行为发起者和行为接受者的潜在角色，避免单向分解丢失入边侧结构。

\section{无标签测试分布特征}

正式测试候选集合本身反映了评测阶段的候选采样机制。方案仅利用公开字段构造转导式无监督特征：

\begin{itemize}
  \item 候选节点在全部测试查询中的出现频次；
  \item 候选节点在测试时间分位窗口中的出现频次；
  \item 候选节点首次/末次出现时间及时间跨度；
  \item 同一源节点与同一候选组合的出现频次、首次/末次时间；
  \item 上述统计的行内归一化值、倒数排名和比例特征。
\end{itemize}

这些特征不包含测试标签。它们的作用是校准训练分布和测试候选生成分布之间的偏移；所有合规边界在第~\ref{sec:compliance} 节进一步说明。

\section{行内相对特征与特征筛选}

原始特征总数为 180。除绝对统计外，对关键候选量生成行内相对变换：

\begin{align*}
\operatorname{norm}(x_{ij})
&=\frac{x_{ij}-\min_j x_{ij}}{\max_j x_{ij}-\min_j x_{ij}+\varepsilon},\\
\operatorname{invRank}(x_{ij})
&=\frac{1}{\operatorname{rank}_{\downarrow}(x_{ij})},\\
\operatorname{zscore}(x_{ij})
&=\frac{x_{ij}-\mu_i}{\sigma_i+\varepsilon}.
\end{align*}

这些变换直接服务于 MRR：即使两个查询的绝对计数尺度不同，模型仍能识别“当前 100 个候选中最强”的节点。

候选在输入行中的列位置不是语义特征，最终方案明确关闭 \texttt{pos\_norm} 和 \texttt{inv\_pos}，避免模型学习偶然的候选排列偏差。dataset1 使用 v65 配置，从 180 维中选择 167 维；dataset2 使用更精简的 v61 配置，选择 105 维，以减轻较稀疏或分布不同的数据上冗余特征带来的过拟合。

\section{排序学习与双分支集成}

\subsection{XGBRanker}

每个查询行是一个长度为 100 的排序 group。标签向量只有真实目标位置为 1，其余为 0。模型使用 XGBRanker，训练目标为 \texttt{rank:ndcg}，评估指标为 \texttt{ndcg@10}。最终共享超参数见表~\ref{tab:xgb}。

\begin{table}[htbp]
\centering
\caption{XGBRanker 共享超参数}
\label{tab:xgb}
\begin{tabular}{ll@{\qquad}ll}
\toprule
参数 & 取值 & 参数 & 取值 \\
\midrule
max\_depth & 5 & learning\_rate & 0.032 \\
subsample & 0.88 & colsample\_bytree & 0.90 \\
min\_child\_weight & 20.0 & reg\_lambda & 4.0 \\
reg\_alpha & 0.0 & random\_state & 20260525 \\
objective & rank:ndcg & eval\_metric & ndcg@10 \\
\bottomrule
\end{tabular}
\end{table}

GPU 可用时使用 CUDA 上的 \texttt{hist}；GPU 失败时自动退回 CPU \texttt{hist}。回退只改变耗时，不改变算法流程和输入输出约定。

\subsection{规则先验与混合}

规则基线 V18 将 pair 频次/时效性、节点热度、最近目标命中、相邻转移、跳跃转移和微小确定性散列噪声线性组合，得到稳定的图启发式分数 $s^{\text{base}}$。它既是独立基线，也以归一化分数、倒数名次、z-score 和距行内最大值差等形式进入排序器特征。

排序器输出与 V18 基线均在每行归一化后再混合。dataset1 保留 5\% 基线，可以降低树模型在极端稀疏查询上的波动；dataset2 的最优设置为纯排序器输出。

\subsection{互补预测分支的排序空间校准}

第二预测分支同样对每行 100 个候选给出分数，但两个分支的绝对尺度不可直接比较。开发阶段先把分数转换为行内名次 $r^{(g)}_{ij}$ 与 $r^{(c)}_{ij}$，再比较下列尺度不敏感的集成形式：

\begin{align*}
s^{\mathrm{RR}}_{ij}
&=\frac{w}{r^{(g)}_{ij}}+\frac{1-w}{r^{(c)}_{ij}},\\
s^{\mathrm{RRF}}_{ij}
&=\frac{w}{k+r^{(g)}_{ij}}+\frac{1-w}{k+r^{(c)}_{ij}},\\
s^{\mathrm{Borda}}_{ij}
&=w\frac{101-r^{(g)}_{ij}}{100}+(1-w)\frac{101-r^{(c)}_{ij}}{100},\\
s^{\mathrm{RP}}_{ij}
&=-\left(w\log r^{(g)}_{ij}+(1-w)\log r^{(c)}_{ij}\right).
\end{align*}

此外比较了行内 min--max 归一化加权和。dataset1 的候选权重主要落在主分支 $0.80$--$0.90$，dataset2 的主分支权重主要落在 $0.15$--$0.25$；这说明两个数据集的优势来源不同，也为最终的数据集级路由提供了证据。

\subsection{置信度门控实验}

为了只在少量强分歧行借用互补信息，定义分支置信度
\[
q_i=\frac{s_{i,(1)}-s_{i,(2)}}{\operatorname{std}(s_i)+\varepsilon},
\]
其中 $s_{i,(1)}$、$s_{i,(2)}$ 为该行前两名分数。对两个分支的 $q_i$ 做稳健中心化和 sigmoid 校准，只在 top-1 不一致的行上，根据置信度差选择约 3\%、5\% 或 8\% 的候选行切换分支。该方法用于验证互补性确实存在，但门控阈值会增加隐藏测试分布漂移风险，因此未作为最终默认方案。

\subsection{最终数据集感知路由}

实验表明，dataset1 应最大程度保留 V69 主分支，而 dataset2 更适合保留互补分支的完整排序。最终 V83 不再为每一行估计复杂权重，而采用确定性路由：

\begin{itemize}
  \item dataset1：选择 V69 的 v65/1200-tree 输出；
  \item dataset2：选择互补预测分支的原始输出；
  \item 打包前对两个选中成员逐行检查 100 列有限浮点数，并核对 SHA256。
\end{itemize}

这种做法的关键不是简单拼接文件，而是先通过排序空间实验识别每个数据集的优势分支，再把分支选择固化为可审计、可重复的推理路由。

\section{优化演进与最终取舍}

整个优化过程围绕“验证分布一致性、图结构表达、过拟合控制、数据集差异化”逐步推进：

\begin{longtable}{C{1.9cm}p{4.2cm}p{8.1cm}}
\caption{主要优化阶段}\label{tab:evolution}\\
\toprule
阶段 & 主要变化 & 得到的认识 \\
\midrule
\endfirsthead
\toprule
阶段 & 主要变化 & 得到的认识 \\
\midrule
\endhead
V3--V18 & 从节点热度逐步加入 pair、时间衰减、最近行为、转移和跳跃转移，并固定为稳定规则基线。 & 推荐信号同时依赖长期关系和短期序列；单一热度或单一 pair 统计不足，V18 可作为后续模型的可靠先验。 \\
V25--V33 & 引入 XGBoost 排序学习，扩展绝对/相对特征；比较含候选位置和不含位置的版本。 & 直接优化行内排序优于固定权重；候选位置可能带来伪相关，最终坚持 no-position。 \\
V34--V43 & 加入长历史、共同邻居、方向 motif、两跳/跳跃关系，并进行特征子集筛选。 & 图特征越多不一定越好；需要按验证表现保留稳定特征，控制稀疏高阶模式的噪声。 \\
V45--V58 & 调整验证行构造、训练组数、树数和 NDCG 目标，使训练候选分布更接近正式测试。 & 从整行候选采样比独立负采样更可靠；模型容量要与验证规模配套。 \\
V59--V61 & 加入 32 维正向 SVD、反向 SVD，并为 dataset2 形成精简的 v61 特征集。 & 低秩表示补足显式统计的泛化能力；有向图需要分别表达入边和出边角色。 \\
V62--V69 & 对 dataset1 采用更丰富的 v65 特征并增至 1200 棵树；dataset2 保持 v61 与 850 棵树；独立确定混合系数。 & 两个数据集的稀疏度和主导信号不同，单一统一配置不如数据集特定配置。V69 成为最终主模型分支。 \\
V71--V82 & 在候选名次空间试验 reciprocal rank、Borda、RRF、rank product、自适应权重和归一化加权。 & 排名变换消除了两个预测分支的分值尺度差异；dataset1 应偏向主分支，dataset2 应偏向互补分支。 \\
V83--V94 & 比较数据集级路由、3\%--8\% 置信度门控以及不同排序校准组合。 & 复杂门控存在阈值漂移风险；V83 直接选择每个数据集的优势分支，结构最简单，成为高分归档版本。 \\
\bottomrule
\end{longtable}

最终取舍的逻辑不是简单“堆更多特征”，而是：先用时间验证保证评价可信，再让候选采样贴近测试分布；随后通过图统计和双向 SVD 补充表达；最后同时在模型容量和预测分支两个层面做数据集特定选择。

\section{最终 V83 配置与结果}

\begin{table}[htbp]
\centering
\caption{主分支配置与最终路由}
\label{tab:final}
\begin{tabularx}{\textwidth}{lC{2.0cm}C{1.55cm}C{1.7cm}C{1.7cm}Y}
\toprule
数据集 & 主分支配置 & 树数 & SVD & 最终来源 & 指标记录 \\
\midrule
dataset1 & v65，167/180 维 & 1200 & 正/反各 32 维 & V69 主分支 & 本地 MRR 0.8456168559 \\
dataset2 & v61，105/180 维 & 850 & 正/反各 32 维 & 互补分支 & 主分支本地 MRR 0.5496001738 \\
主分支合计 & V69 & --- & --- & --- & 本地 MRR 1.3952170297 \\
最终提交 & V83 & --- & --- & 数据集感知路由 & \textbf{A 榜 1.4591} \\
\bottomrule
\end{tabularx}
\end{table}

这里的 V69 本地 MRR 来自固定随机种子和代码所述时间验证过程，用于复现主分支模型选择逻辑；A 榜 1.4591 才是最终 V83 的提交成绩。两者来自不同数据，不能把 1.3952170297 写成最终榜单分数。

\section{工程实现与资源优化}

\subsection{紧凑索引}

pair 和 transition 使用整数复合键
\[
\operatorname{key}(a,b)=a\cdot B+b,
\]
其中 $B$ 大于最大节点编号。复合键排序后以 \texttt{searchsorted} 批量查询，避免维护大量 Python 字典对象。密集节点统计使用 NumPy 数组，特征矩阵统一为 \texttt{float32}。

\subsection{分块构建与预测}

训练和预测按查询行分块生成特征，避免一次性展开全部 $N\times100\times F$ 数据。默认排序器块大小为 1500 行；内存不足时可进一步减小。输出也分块写入 CSV，因此峰值内存主要由历史索引、SVD 和单个特征块决定。

\subsection{确定性与失败回退}

主随机种子固定为 20260525，NumPy 采样、SVD 和 XGBoost 均使用确定的种子。程序在 GPU 训练异常时重建 CPU 模型继续训练；只有 dataset1 和 dataset2 两个结果文件均存在且通过全量矩阵检查时才生成 \texttt{result.zip}，降低半成品被误提交的风险。高分归档中 dataset1 成员 SHA256 为 \texttt{829b6877...cdb7ee5}，dataset2 成员 SHA256 为 \texttt{a56c257f...f044fb4a}，最终 ZIP SHA256 为 \texttt{5b9e61fd...cdc7e6b}；完整值见 \path{records/final_metrics.json}。

\section{复现环境}

审核环境固定为：Ubuntu 22.04、CUDA 12.4、Python 3.10、Jittor 1.3.10.0。候选排序器实际使用 XGBoost 2.1.4；Jittor 版本按比赛审核环境固定，完整版本见 \path{requirements.txt} 和 \path{environment.yml}。

\subsection{安装命令}

\begin{lstlisting}[language=bash]
sudo apt-get update
sudo apt-get install -y python3.10-dev g++ build-essential libomp-dev
python3.10 -m venv .venv
source .venv/bin/activate
python -m pip install --upgrade pip
python -m pip install -r requirements.txt
python code/check_environment.py
\end{lstlisting}

\subsection{数据目录}

\begin{lstlisting}
data/track1/
|-- dataset1/train.csv
|-- dataset1/test.csv
|-- dataset2/train.csv
`-- dataset2/test.csv
\end{lstlisting}

也可以提前把官方压缩包放在 \path{data/track1_data.zip}。代码会在数据目录缺失时从脚本中固定的主办方地址下载并解压。

逐字节复现最终 V83 还需将归档的互补预测分支放为 \path{data/complementary_result.zip}。其 SHA256 必须为
\begin{lstlisting}
974582aa2f8f156bbe7f539487450e3e1d7ca75cfaec6027ee2999ad53fdb461
\end{lstlisting}
该分支作为固定模型输出输入；集成程序校验其文件指纹和全部预测矩阵，但不把其内部训练过程声明为 \path{code/main.py} 的产物。

\subsection{一键运行}

\begin{lstlisting}[language=bash]
chmod +x code/run.sh
chmod +x code/run_primary.sh
bash code/run.sh
\end{lstlisting}

脚本先调用 \path{code/run_primary.sh} 生成 \path{outputs/primary_v69.zip}，再执行：

\begin{lstlisting}[language=bash]
python code/ensemble.py \
  --primary_result outputs/primary_v69.zip \
  --complementary_result data/complementary_result.zip \
  --output_dir outputs/high_score \
  --result_zip outputs/result.zip
\end{lstlisting}

若内存紧张，可执行：

\begin{lstlisting}[language=bash]
bash code/run.sh --ranker_chunk_rows 800 --chunk_rows 6000
\end{lstlisting}

\subsection{预期输出与检查}

\begin{lstlisting}
outputs/
|-- primary_v69.zip
|-- high_score/
|   |-- dataset1.csv       # 61,051 rows, 100 scores per row
|   |-- dataset2.csv       # 153,420 rows, 100 scores per row
|   `-- ensemble_metadata.json
`-- result.zip
\end{lstlisting}

\begin{lstlisting}[language=bash]
unzip -l outputs/result.zip
wc -l outputs/high_score/dataset1.csv outputs/high_score/dataset2.csv
\end{lstlisting}

每行应恰好包含 100 个有限浮点分数；行内最小值约为 0、最大值约为 1；压缩包根目录只包含 \texttt{dataset1.csv} 和 \texttt{dataset2.csv}。

历史 V83 ZIP 的整包 SHA256 记录了当时的容器时间戳。重新打包后容器 SHA256 可以变化；预测版本以两个 CSV 成员的 SHA256 为准，\path{code/ensemble.py} 默认强制核对完整成员指纹。

\section{合规与复现边界}
\label{sec:compliance}

\begin{itemize}
  \item 训练只使用 \texttt{train.csv} 中公开的历史边。
  \item 测试文件只使用 \texttt{src}、\texttt{time} 和 \texttt{c1--c100} 候选集合。
  \item 测试候选频次、时间窗口和源--候选共现均为无标签统计，不包含真实目标。
  \item 验证真实目标来自训练边的严格时间后移切分，不使用未来边构建历史特征。
  \item V69 主分支可从官方数据完整重建，代码不访问隐藏测试标签。
  \item V83 的第二预测分支以固定模型输出归档；\path{code/ensemble.py} 校验其 SHA256、ZIP CRC、行数、列数和有限值，并复现最终推理路由。
  \item 文档不把第二分支的内部训练过程误写为主模型实现；可验证范围是主分支完整训练，以及最终分支选择、校验和打包。
  \item 本地 V69 MRR 与最终 V83 榜单成绩分开报告，避免混淆验证集指标与隐藏测试指标。
\end{itemize}

\section{代码结构与算法对应关系}

\begin{tabularx}{\textwidth}{p{4.0cm}Y}
\toprule
文件 & 作用 \\
\midrule
\path{code/main.py} & 数据下载与读取、时间验证、历史索引、双向 SVD、特征构建、XGBRanker 训练、分块预测及 V69 主分支打包。 \\
\path{code/ensemble.py} & V83 数据集感知双分支路由、全量矩阵校验、成员指纹记录和最终 ZIP 打包。 \\
\path{code/config.json} & V69 主分支参数、V83 路由、榜单成绩与产物指纹。 \\
\path{code/run.sh} & Ubuntu 环境下最终高分版本一键入口。 \\
\path{code/run_primary.sh} & 只复现 V69 主分支的入口。 \\
\path{code/check_environment.py} & Python、Jittor、XGBoost 等依赖版本检查。 \\
\path{records/final_metrics.json} & 最终指标、特征数与输出行数记录。 \\
\path{records/artifact_audit.md} & 高分 ZIP、成员来源、哈希、行列数及其他候选版本的审计记录。 \\
\path{records/reproduction.log} & V69 dataset1 完整训练日志；末尾打包错误来自隔离运行未同时生成 dataset2。 \\
\path{docs/复现说明.md} & 精简的命令行复现与故障排查说明。 \\
\bottomrule
\end{tabularx}

\section{常见问题}

\begin{itemize}
  \item \textbf{GPU 训练失败：}代码会自动退回 CPU \texttt{hist}，结果逻辑不变，但运行时间明显增加。
  \item \textbf{数据下载失败：}手动将官方数据保存为 \path{data/track1_data.zip}，或直接解压为规定目录结构。
  \item \textbf{内存不足：}降低 \texttt{--ranker\_chunk\_rows}；完整复现仍建议至少 32 GB 内存。
  \item \textbf{互补分支缺失：}将归档文件放到 \path{data/complementary_result.zip}，或设置环境变量 \texttt{COMPLEMENTARY\_RESULT\_ZIP}。
  \item \textbf{文件指纹不一致：}停止提交并核对输入归档；可在 \path{records/final_metrics.json} 查询完整 SHA256。
  \item \textbf{输出 ZIP 未生成：}根据 \path{code/ensemble.py} 的错误信息检查 ZIP 成员、行数、列数和非法浮点数。
  \item \textbf{依赖版本告警：}优先使用仓库中固定的 \path{requirements.txt}，尤其是 XGBoost 2.1.4。
\end{itemize}

\section{结论}

最终方案的核心并非某一个单独模型，而是五层逻辑的闭环：严格时间验证保证优化方向可信；候选行模拟使本地训练和正式评测一致；显式图统计与双向低秩表示共同覆盖记忆和泛化；数据集特定的特征筛选与容量控制减少过拟合；最后通过排序空间校准实验识别不同数据集的优势预测分支，并固化为确定性路由。V69 是可完整训练复现的主分支，V83 是经过文件指纹与输出结构审计的最终高分版本，A 榜成绩为 1.4591。

\end{document}
